\section{Experiments: Is Agent Creativity Axis-Dependent?}
\label{sec:experiments}

\paragraph{What was run.} One run: $50$ challenges $\times$ $5$ entrants $=250$ records, all scored,
under live provider-native web search, judged by the three-vendor full panel of \S\ref{sec:judge}.
The entrants are \texttt{openai\_deepresearch} and \texttt{openai\_shallow} (OpenAI),
\texttt{gemini\_deepresearch\_max} and \texttt{gemini\_shallow} (Google), and
\texttt{claude\_shallow} (Anthropic). There is \textbf{no} Perplexity entrant and \textbf{no}
open-weight entrant; earlier drafts named both, and neither exists in the released registry. The
numbers below are promoted into the tracked \texttt{result/leaderboard.json}; recomputing them
requires the run's \texttt{records.jsonl}, which is \emph{not} distributed (see the reproducibility
paragraph at the end of this section).

\paragraph{Stated prior expectations, not a pre-registration.} Per-category directions are recorded
in \texttt{analysis/preregistration.yaml}: under a mechanistic account, agents coupling relational
grounding with convergent refinement should excel at cross-domain (C3, structure-mapping
\citep{gentner1983structure}) and counterfactual (C2 \citep{byrne2005rational}) synthesis, be at
parity on meta-creativity (C4), and show nulls on constrained (C1) and open-ended (C5) synthesis
\citep{boden2004creative,campbell1960blind}. This is \textbf{not} an externally timestamped
pre-registration, and the file says so: it was created in-repo after the code, so it cannot evidence
ordering. Treat the directions as the authors' stated prior expectation.

\begin{table}[t]
\centering
\small
\begin{tabular}{lccccccccc}
\toprule
 & \multicolumn{1}{c}{Primary} & \multicolumn{3}{c}{Per-axis mean} & \multicolumn{5}{c}{Primary endpoint by category} \\
\cmidrule(lr){2-2}\cmidrule(lr){3-5}\cmidrule(lr){6-10}
Entrant & $\min(N,U,S)$ & $N$ & $U$ & $S$ & C1 & C2 & C3 & C4 & C5 \\
\midrule
\texttt{openai\_deepresearch}     & 6.24 & 8.52 & 6.76 & 7.90 & 4.40 & 7.70 & 6.70 & 6.00 & 6.40 \\
\texttt{openai\_shallow}          & 6.18 & 8.58 & 6.68 & 7.82 & 3.80 & 7.60 & \textbf{7.50} & 5.90 & 6.10 \\
\texttt{claude\_shallow}          & 5.72 & 8.08 & 6.06 & 7.44 & 3.20 & 7.30 & 6.30 & 5.80 & 6.00 \\
\texttt{gemini\_deepresearch\_max}& 4.92 & 7.96 & 5.02 & 7.42 & 3.60 & 5.80 & 5.50 & 4.30 & 5.40 \\
\texttt{gemini\_shallow}          & 4.18 & 7.56 & 4.34 & 6.96 & 3.40 & 4.60 & 4.70 & 4.60 & 3.60 \\
\midrule
Field mean                        & 5.45 & 8.14 & 5.77 & 7.51 & 3.68 & 6.60 & 6.14 & 5.32 & 5.50 \\
\bottomrule
\end{tabular}
\caption{CREST results, $n=50$ per entrant ($250$ records, all scored). \textbf{These are not
rankings.} The between-entrant spread on the primary endpoint is $2.06$ points, while judge severity
spans $3.32/3.20/3.74$ points per axis (\S\ref{sec:judge}), so the ordering is not identifiable from
this run. Two readings that \emph{are} supported: (i) \textbf{Useful is the binding axis}---per-axis
means run $7.56$--$8.58$ for New but only $4.34$--$6.76$ for Useful, so a conjunctive endpoint on
these items is largely measuring usefulness; (ii) \textbf{C1 is uniformly hard}, at $3.20$--$4.40$
for every entrant, the only category where no entrant clears $5$.}
\label{tab:results}
\end{table}

\paragraph{Construct threat: the primary endpoint is largely hard-rule compliance.} The conjunctive endpoint is dominated by one axis. On the released run $\mathrm{corr}(\min(N,U,S), U) = 0.976$, the composite \emph{equals} Useful in \textbf{190 of 250 records} ($76\%$), and \emph{New} is the unique minimum in \textbf{1 of 250}. Because \emph{Useful} is defined so that a violated hard rule caps it at ${\leq}3$, the endpoint behaves close to a rule-violation indicator, and the between-category structure tracks constraint density rather than creative process: category means order C1 $3.68 <$ C4 $5.32 <$ C5 $5.50 <$ C3 $6.14 <$ C2 $6.60$, in near-perfect correspondence with the share of judge rows citing a rule violation. Compounding this, the only axis with any one-way reliability is Useful ($\mathrm{ICC}(1,k)=0.618$); the two creativity-bearing axes sit at $0.0$, and \emph{New} averages $8.14/10$---near ceiling and barely discriminating even \emph{with} the prior-art list supplied. We therefore cannot claim the endpoint measures creativity rather than constrained-generation compliance, and the C1 ``predicted null'' is indistinguishable from a floor effect on the items carrying the most checkable rules. This is the most serious open threat to the construct and it is not resolved here.

\paragraph{Axis dependence is visible; the crossover is weak.} Task type moves scores far more than
entrant identity does: the field mean ranges from $3.68$ (C1) to $6.60$ (C2), a $2.92$-point spread,
against $2.06$ points between the best and worst entrant. Within that, one rank reversal survives:
\texttt{openai\_shallow} leads the field on cross-domain C3 ($7.50$ vs $6.70$) while
\texttt{openai\_deepresearch} leads overall and on C2, C4 and C5. That is a genuine
agent$\times$category interaction, but it is \emph{one} reversal between two entrants of the same
vendor, separated by $0.06$ points overall---far weaker than the ``no single creativity scalar
orders current agents'' claim earlier drafts made, and we withdraw that claim. With severity
exceeding the between-entrant spread, we cannot distinguish a real reversal from judge noise.

\paragraph{The predicted pattern is not confirmed.} The stated expectation was wins on C2/C3, parity
on C4, nulls on C1/C5. What the data show is a field-wide ordering by task difficulty
(C2 $>$ C3 $>$ C5 $>$ C4 $>$ C1) that is similar for every entrant, which is a property of the
\emph{items} rather than a dissociation among agents. The C1 null holds in the weak sense that no
entrant does well there, but that is consistent with C1 simply being harder, not with a mechanism
being absent.

\paragraph{Analyses claimed by earlier drafts that are NOT performed.} Named here rather than
omitted, because each was previously asserted as done. There is \textbf{no} rater random effect:
\texttt{analysis/stats.py} fits a challenge-level random intercept only, so the ``crossed random
effects for challenge and rater'' of earlier drafts is withdrawn. There are \textbf{no} MDE or power
curves, \textbf{no} multiplicity correction across the reported cells, \textbf{no} output-length
covariate, and \textbf{no} judge-free diversity check. Budgets are \textbf{not} matched: realised
search counts run from $2.9$ (\texttt{openai\_shallow}) to ${\sim}76$
(\texttt{gemini\_deepresearch\_max}), so depth is a factor, not noise, and no cost--creativity
frontier is reported. The registered TOST is also \textbf{not usable as registered}: it tests a
category composite against a reference of $0$ with equivalence bounds of $\pm1$ on a $1$--$10$
scale, which no category can satisfy, so the null categories are reported as descriptive means
rather than as equivalence results.

\paragraph{What a reader can and cannot reproduce.} The item set, the harness, the judge
configuration, the analysis code and the promoted results are all tracked. The $250$-record run
directory is \textbf{not}: \texttt{result/runs/} is gitignored, so per-record judge rows and agent
transcripts are not in the repository and the table above cannot be recomputed from a fresh clone.
Re-running the benchmark re-queries the live web, so it will not reproduce these numbers exactly
under any circumstances. A frozen-index mode that would make the run reproducible is future work and
\texttt{harness/agents.py::FrozenIndex} deliberately raises rather than pretending to provide it.
